\documentclass[11pt]{article}
\usepackage[margin=1in]{geometry}
\usepackage{amsmath, amssymb, amsthm}
\usepackage{booktabs}
\usepackage{enumitem}
\usepackage[hidelinks]{hyperref}

\newcommand{\cop}{\mathrm{cop}}
\newcommand{\pcf}{\mathrm{pcf}}
\newcommand{\code}[1]{\texttt{#1}}
\newcommand{\SVD}{\mathrm{SVD}}
\DeclareMathOperator*{\argmin}{arg\,min}

\title{EM completion of the factor loadings, and copula-specific factors}
\author{Implementation notes for \code{pca\_method = :em} and \code{n\_copula\_factors}}
\date{September 14, 2026}

\begin{document}
\maketitle

\section{What these two options do}

Both change how the loading matrix $\Lambda$ --- the map from factors to
coefficients --- is estimated. Neither touches the state-space recursions, the
priors, or the likelihood: everything downstream reads the factor count off
$\mathrm{size}(\Lambda, 2)$ and is otherwise unchanged.

\begin{itemize}[leftmargin=*]
  \item \code{pca\_method = :em} estimates $\Lambda$ from \emph{every}
        dataset-period, using each one for the coefficients it actually
        observes, instead of discarding any period where some measure is
        missing.
  \item \code{n\_copula\_factors = $r_c$} adds $r_c$ factors that only the
        copula coefficients load on, so the dependence structure has somewhere
        to move that is not already spoken for by the marginals.
\end{itemize}

Motivation and the diagnosis that led here are in
\code{doc/copula\_identification.tex}; this note is the math of what the code
does.

\section{The pool and its missingness}

Stack the demeaned, object-scaled coefficient vectors of all datasets into
\begin{equation}
  P \in \mathbb{R}^{n \times N},
  \qquad
  n = n_c + n_q, \qquad N = \sum_j T_j ,
\end{equation}
where each \emph{column} is one (dataset, period) observation vector, each
\emph{row} one coefficient (the first $n_c$ rows are copula coefficients, the
remaining $n_q = D\cdot G$ the marginal percentile functions). Let
\begin{equation}
  \Omega = \bigl\{ (i,t) : P_{it} \text{ observed} \bigr\}
\end{equation}
and write $\mathcal{P}_\Omega$ for the projection that keeps entries in $\Omega$
and zeroes the rest. The missingness is \textbf{block structured}: a dataset
that does not measure, say, consumption has \emph{all} consumption-related rows
missing in \emph{all} of its columns.

\paragraph{Baseline (\code{:complete\_cases}).} The code retains only columns
with no missing entry,
\begin{equation}
  \mathcal{K} = \bigl\{ t : (i,t) \in \Omega \ \ \forall i \bigr\},
  \qquad
  \hat\Lambda_{\mathrm{cc}} = \SVD_r\bigl( P_{\cdot\mathcal{K}} \bigr),
\end{equation}
$\SVD_r$ denoting the leading $r$ left singular vectors (scaled by the square
roots of the eigenvalues, as in the code). Only a dataset observing \emph{all}
measures contributes, at its wave dates: in both the replica exercise and the
empirical model $|\mathcal{K}| = 12$, so $\mathrm{rank}(P_{\cdot\mathcal{K}})
\le 11$ and each row's $r$ loadings are fitted with $12 - r$ degrees of freedom.

\section{EM completion}

\subsection{The objective}

Fit the best rank-$r$ matrix to the \emph{observed} entries only:
\begin{equation}
  \hat Z \;=\; \argmin_{\mathrm{rank}(Z)\le r}\;
  \bigl\| \mathcal{P}_\Omega\bigl(P - Z\bigr) \bigr\|_F^2 .
  \label{eq:objective}
\end{equation}
This is the Gaussian PPCA likelihood in the zero-noise limit; with isotropic
noise $\sigma^2 > 0$ the E-step would additionally shrink the imputations
toward zero by the factor $\sigma^2/(\sigma^2 + s_k^2)$ per singular direction.
The implementation uses the zero-noise version, which is the standard
impute-and-refit EM (also known as Hard-Impute).

\subsection{The iteration}

Initialize $X^{(0)} = \mathcal{P}_\Omega(P)$, i.e.\ missing entries at $0$ ---
the unconditional mean, since the pool is demeaned. Then iterate
\begin{align}
  \text{M-step:}&\quad Z^{(s)} = \SVD_r\bigl( X^{(s)} \bigr), \label{eq:mstep}\\
  \text{E-step:}&\quad X^{(s+1)} = \mathcal{P}_\Omega(P) \;+\; \mathcal{P}_\Omega^{\perp}\bigl(Z^{(s)}\bigr),
  \label{eq:estep}
\end{align}
so observed entries are held at their data values and missing entries are
replaced by their current rank-$r$ reconstruction. Each half-step is a
minimization of \eqref{eq:objective} in one argument, so the objective is
non-increasing and the iteration converges to a stationary point. The code stops
when the root-mean-square change over the imputed entries,
\begin{equation}
  \Delta^{(s)} = \Bigl( |\Omega^{\perp}|^{-1} \textstyle\sum_{(i,t)\notin\Omega}
  \bigl(Z^{(s)}_{it} - X^{(s)}_{it}\bigr)^2 \Bigr)^{1/2},
\end{equation}
falls below $10^{-7}$ or stops changing, with a cap of $200$ iterations.
Columns observing less than a fraction \code{min\_frac} $= 0.05$ of the
coefficients are dropped before the iteration; they carry almost no information
and would otherwise be filled in almost entirely by the model.

The loadings are then
\begin{equation}
  \hat\Lambda_{\mathrm{em}} = \SVD_r\bigl( \hat X \bigr),
  \qquad \hat X = \lim_s X^{(s)} .
\end{equation}

\subsection{Loadings from everything, scores from the complete cases}

Only $\Lambda$ is taken from the completed matrix. The factor scores --- which
feed the Minnesota prior's scaling and several downstream objects --- are still
computed on the complete-case columns,
\begin{equation}
  \hat F_{\mathcal{K}} = \hat\Lambda_{\mathrm{em}}^{\top} X_{\cdot\mathcal{K}} ,
\end{equation}
so every array downstream keeps its existing shape and the change is isolated to
the loading matrix. This is deliberate: the completed columns are
dataset-blocks concatenated in time, not a single clean time series, so using
them as a score path would silently feed a spurious ordering into anything that
treats the score matrix dynamically.

\subsection{What EM can and cannot recover}

For a row $i$ observed on the set $\mathcal{T}_i$, the loading estimate behaves
like the least-squares fit
\begin{equation}
  \hat\lambda_i \approx \Bigl( \sum_{t\in\mathcal{T}_i} \hat f_t \hat f_t^{\top} \Bigr)^{-1}
                        \sum_{t\in\mathcal{T}_i} \hat f_t\, P_{it},
\end{equation}
an $r$-regressor problem with $|\mathcal{T}_i|$ observations, rather than a
joint decomposition of a twelve-column matrix. Consistency requires
$|\mathcal{T}_i| \to \infty$ and the factors to be estimable from the rows that
are widely observed.

The block structure sets a hard limit. If two groups of coefficients are never
observed in the same dataset-period, no completion scheme creates information
about their joint variation: the imputations in those cells are generated by
the model, so they can propagate the factor structure but cannot test it. In our
pools, consumption and wealth coefficients co-occur only in the
PSID-type dataset, so cross-measure dependence still rests on those dates. What
EM recovers is the within-measure and two-dimensional-copula information from
every \emph{other} dataset-period, which the complete-case rule discards
entirely.

\section{Copula-specific factors}

\subsection{Extraction}

Let $\hat\Lambda$ be the common loadings from either method and $\hat F$ the
corresponding scores on the columns those loadings were estimated on. The
residual
\begin{equation}
  R = X - \hat\Lambda \hat F
\end{equation}
is what the common factors leave unexplained. Restrict to the copula rows and
take its leading $r_c$ components,
\begin{equation}
  R_c = R_{1:n_c,\,\cdot} \in \mathbb{R}^{n_c \times N_{\mathrm{load}}},
  \qquad
  \Lambda_{cc} = \SVD_{r_c}\bigl( R_c \bigr) \in \mathbb{R}^{n_c \times r_c} .
\end{equation}
Under \code{:em} these loadings are estimated on the completed columns, the same
ones the common loadings used; the scores are again taken on the complete-case
columns so the score matrix stays conformable.

\subsection{The augmented loading matrix}\label{sec:blocktri}

Partitioning $\hat\Lambda$ into its marginal and copula rows,
$\hat\Lambda = \bigl(\Lambda_{qm}^{\top},\, \Lambda_{cm}^{\top}\bigr)^{\top}$,
the model's loading matrix becomes block-triangular:
\begin{equation}
  \begin{pmatrix} \theta^{\pcf}_{t} \\[2pt] \theta^{\cop}_{t} \end{pmatrix}
  \;=\;
  \underbrace{\begin{pmatrix} \Lambda_{qm} & 0 \\ \Lambda_{cm} & \Lambda_{cc} \end{pmatrix}}_{\textstyle \Lambda_{\mathrm{aug}}}
  \begin{pmatrix} f^m_t \\ f^c_t \end{pmatrix} + e_t ,
  \qquad
  f_t = \begin{pmatrix} f^m_t \\ f^c_t \end{pmatrix} \in \mathbb{R}^{r + r_c}.
  \label{eq:blockload}
\end{equation}
The zero block is structural, not a prior: marginal coefficients load exactly
zero on the copula factors, so those factors cannot distort the marginals. The
state equation is unchanged in form,
\begin{equation}
  \begin{pmatrix} f_t \\ a_t \end{pmatrix} =
  \begin{pmatrix} A & B \\ C & D \end{pmatrix}
  \begin{pmatrix} f_{t-1} \\ a_{t-1} \end{pmatrix} + \eta_t ,
\end{equation}
now with $A \in \mathbb{R}^{(r+r_c)\times(r+r_c)}$ --- so the copula factors may
respond to the aggregates through $B$ and feed back through $C$, which is what
lets dependence move over the business cycle. Under
\code{persistence\_prior = :per\_factor} each copula factor also receives its
own persistence target $\rho_{F,j}$.

\subsection{Identification}

Two indeterminacies must be pinned down. The zero block in
\eqref{eq:blockload} restricts admissible rotations $\Lambda_{\mathrm{aug}}
\mapsto \Lambda_{\mathrm{aug}} H$ to those preserving it, i.e.\ block-triangular
$H$; combined with the PCA normalization ($\Lambda^{\top}\Lambda$ diagonal with
ordered entries, within each block) the loadings are fixed up to sign. Sign
conventions follow whatever \code{MultivariateStats} returns and are irrelevant
to the likelihood.

\paragraph{An honest caveat.} The copula factors are extracted from the residual
of the common fit, so their \emph{scores} are close to uncorrelated with the
common scores --- but the \emph{loading vectors} are not orthogonal. The
residual satisfies $\hat\Lambda^{\top} R = 0$, a sum over \emph{all} rows;
restricting to copula rows breaks it, and in general
\begin{equation}
  \hat\Lambda^{\top} \begin{pmatrix} 0 \\ \Lambda_{cc} \end{pmatrix}
  \;=\; \Lambda_{cm}^{\top}\Lambda_{cc} \;\neq\; 0 .
\end{equation}
The state space tolerates this --- $\Omega$ is unrestricted and the VAR allows
interaction, so correlated factors are a reparameterization, not a
misspecification --- but the label ``copula-specific'' is approximate: $f^c$
carries a component that the common factors could also have represented. If a
cleaner decomposition is wanted, orthogonalizing the copula scores against
$\hat F$ (Gram--Schmidt) before appending would enforce it at the cost of
changing $\Lambda_{cc}$.

\subsection{Not yet implemented: shrinkage on the copula factors}

The design intent is that $r_c$ can be set generously and switched off by the
data. That requires a global--local prior on the copula factors' innovation
scales, in the same form as the trend block's,
\begin{equation}
  \omega_{c,\ell} \sim N^{+}\bigl(0, (\lambda_c \psi_\ell)^2\bigr), \qquad
  \psi_\ell \sim C^{+}(0,1), \qquad
  \lambda_c \sim C^{+}(0, s_c), \qquad \ell = 1,\ldots,r_c .
\end{equation}
Until that is added, the copula factors use the standard $\Omega_f$ prior, so
$r_c$ should be chosen conservatively --- there is nothing pushing them to zero
when they are not needed.


\section{Three loading structures, and where information flows}
\label{sec:structures}

Section \ref{sec:blocktri} added copula-specific factors on top of a shared
block. The natural-looking alternative is to give each block its own factors
outright. The two are not close substitutes, and the reason is about
\emph{information}, not flexibility.

Write the measurement equation for one dataset as $y_t = G x_t + e_t$ with
$e_t \sim N(0,\Sigma)$, $\Sigma$ diagonal, and split the coefficient rows into
the copula set $\mathcal{C}$ ($n_c \approx 1694$ rows) and the marginal set
$\mathcal{M}$ ($n_q = 36$ rows). The three structures are
\begin{equation}
  \Lambda_{\text{joint}} = \begin{pmatrix} \Lambda_{qm} \\ \Lambda_{cm} \end{pmatrix},
  \qquad
  \Lambda_{\triangle} = \begin{pmatrix} \Lambda_{qm} & 0 \\ \Lambda_{cm} & \Lambda_{cc}\end{pmatrix},
  \qquad
  \Lambda_{\mathrm{diag}} = \begin{pmatrix} \Lambda_{qm} & 0 \\ 0 & \Lambda_{cc}\end{pmatrix},
\end{equation}
with $f_t = (f^m_t, f^c_t)$ in the latter two.

\subsection{How much each block informs the states}

The filter's update adds to the state precision
\begin{equation}
  G^{\top}\Sigma^{-1}G \;=\; \sum_{i} \frac{g_i g_i^{\top}}{\sigma_i^2}
  \;=\; \underbrace{\sum_{i\in\mathcal{C}} \frac{g_i g_i^{\top}}{\sigma_c^2}}_{\text{copula rows}}
  \;+\; \underbrace{\sum_{i\in\mathcal{M}} \frac{g_i g_i^{\top}}{\sigma_m^2}}_{\text{marginal rows}},
  \label{eq:precision}
\end{equation}
$g_i$ being row $i$ of $G$ (the loading row, premultiplied). Restrict
\eqref{eq:precision} to the marginal factors' block:
\begin{align}
  \text{joint / triangular:}&\quad
  \bigl[G^{\top}\Sigma^{-1}G\bigr]_{mm} =
  \sum_{i\in\mathcal{M}} \frac{\lambda^{qm}_i \lambda^{qm\top}_i}{\sigma_m^2}
  \;+\; \sum_{i\in\mathcal{C}} \frac{\lambda^{cm}_i \lambda^{cm\top}_i}{\sigma_c^2},
  \label{eq:info-tri}\\[2pt]
  \text{block-diagonal:}&\quad
  \bigl[G^{\top}\Sigma^{-1}G\bigr]_{mm} =
  \sum_{i\in\mathcal{M}} \frac{\lambda^{qm}_i \lambda^{qm\top}_i}{\sigma_m^2}
  \qquad\text{only.}
  \label{eq:info-diag}
\end{align}
The second term of \eqref{eq:info-tri} is the whole point: under the joint and
triangular structures the $1694$ copula rows also pin down the factors the
marginals load on, because those rows carry $\Lambda_{cm} \neq 0$. Setting that
block to zero removes it. With the estimated scales at the mode
($\sigma^2_{a,\cop} \approx 37$, $\sigma^2_{a,\mathrm{inc}} \approx 8.5$) the
direct information about the marginal factors falls by a factor of roughly
\begin{equation}
  \frac{n_c/\sigma_c^2}{n_q/\sigma_m^2}
  \;\approx\; \frac{1694/37}{36/8.5} \;\approx\; 11 ,
\end{equation}
i.e.\ about an order of magnitude. The marginals had been riding on the
copula's measurement base, and block-diagonal loadings take that away.

\subsection{The only channel that remains, and the prior that closes it}

Under $\Lambda_{\mathrm{diag}}$ the copula data can still reach the marginal
states, but only \emph{indirectly}, through the prediction step. The update
\begin{equation}
  \mathbb{E}\bigl[f^m_t \mid y_{1:t}\bigr]
  = \hat f^m_{t|t-1}
  + P_{t|t-1}^{mc}\, G_c^{\top}\bigl(G_c P_{t|t-1} G_c^{\top} + \Sigma_c\bigr)^{-1}
    \bigl(y^{\cop}_t - \hat y^{\cop}_t\bigr) + \dots
  \label{eq:transmission}
\end{equation}
transmits a copula surprise to $f^m$ in proportion to the cross-block state
covariance $P^{mc}_{t|t-1}$, which is generated only by the off-diagonal blocks
of $\Omega$ and of $A$. If $\Omega_{mc} = 0$ and $A$ is block-diagonal, then
$P^{mc} = 0$ and the transmission is exactly zero.

That channel is therefore load-bearing --- and it is the one the prior shrinks.
The innovation correlation matrix carries
$\mathrm{LKJ}(\nu)$ with $\nu = 1 + \log(1+e^{h})$ estimated, and $\nu > 1$
tilts toward the identity, i.e.\ toward $\Omega_{mc} = 0$. Under joint or
triangular loadings comovement between the blocks is free, being built into the
loadings; under block-diagonal it must be bought against a prior that leans the
other way.

\subsection{What happened when this was run}

\code{loading\_structure = :block} with four marginal and four copula factors,
on top of EM loadings and untransformed aggregates:

\begin{center}
\begin{tabular}{lrrr}
\toprule
$R^2$ to truth & baseline & EM + levels & block $4\times4$ \\
\midrule
wealth, top quintile (a)      & 57 & 58 & $-28$ \\
wealth, middle-40 share (a)   &  7 & 28 & $-143$ \\
wealth, top-10 share (d)      & 57 & 65 & $-47$ \\
consumption, middle q3 (a)    & 74 & 67 & 19 \\
income, bottom quintile (e)   & 99 & 87 & 98 \\
income, top-10 share (a)      & 93 & 84 & 89 \\
\bottomrule
\end{tabular}
\end{center}

Income --- the measure observed by every dataset, and the one least dependent on
borrowed information --- returned to its baseline fit, which is what the
structure was designed to protect. Everything else collapsed. The mode shows the
mechanism directly: with the factor order (copula $1$--$4$, marginal $5$--$8$),
the estimated innovation scales were
\begin{equation}
  \omega_f = (\,0.02,\; 1.12,\; 2.66,\; 0.42 \;\mid\; 0.69,\; 0.02,\; 0.00,\; 0.01\,),
\end{equation}
so three of the four marginal factors were estimated at zero: they could not be
identified from $36$ rows once \eqref{eq:info-tri}'s second term was removed.
Their trend states absorbed the movement instead
($\omega_\tau = 0.069,\,0.274,\,0.142$ on those slots), and random walks of that
size are what produced the wandering paths and isolated spikes in the panels.

\subsection{Conclusion}

Block-\emph{triangular} dominates block-diagonal for this model. It keeps
$\Lambda_{cm} \neq 0$, so the copula's $1694$ rows continue to identify the
factors that reconstruct the marginals, while $\Lambda_{cc}$ gives the copula
private directions the marginals cannot be distorted by. Block-diagonal buys
symmetry at the price of the term in \eqref{eq:info-tri} that was doing the
work.

A caveat on how the two counts were chosen. The cross-validation in
\code{doc/em\_rank\_selection.tex} put the marginal rows' own rank at four. That
is a statement about \emph{reconstruction given the whole panel}, not about
\emph{identification inside the filter from 36 rows}; this run is what the
difference between those two questions looks like in practice.

\section{Dimensionality}

\code{:em} adds no parameters: it changes how $\Lambda$ is estimated, not what
is estimated. Copula factors raise the distributional factor count from $r$ to
$r + r_c$, and the cost is dominated by the innovation correlation block, which
grows quadratically. For the current replica spec ($q = 8$ aggregate factors,
per-factor persistence, stochastic trend):

\begin{center}
\begin{tabular}{ccrrrrrr}
\toprule
$r_c$ & $r + r_c$ & total & $A,B,C,D$ & $\Omega$ & $\Sigma$ & $H$ & $\tau$ \\
\midrule
0 & 5 & 247 & 113 & \phantom{0}91 & 22 & 10 & 11 \\
1 & 6 & 291 & 140 & 105 & 22 & 11 & 13 \\
2 & 7 & 338 & 169 & 120 & 22 & 12 & 15 \\
3 & 8 & 388 & 200 & 136 & 22 & 13 & 17 \\
\bottomrule
\end{tabular}
\end{center}

\section{What the run log reports}

\begin{itemize}[leftmargin=*]
  \item \code{Distributional PCA: EM completion --- loadings from $N$
        dataset-periods across $k$ dataset(s) (complete-case columns: 12); $s$
        iterations, final change $\Delta$}. The jump from 12 to $N$ is the point
        of the exercise. The spectrum and factor-count criteria printed nearby
        still describe the complete-case pool: the completed matrix is rank-$r$
        by construction, so its own spectrum is degenerate and would be
        misleading.
  \item \code{Copula-specific factors: added $r_c$ (distributional factor count
        $r \to r + r_c$); loadings from $N$ columns; they explain $X\%$ of the
        copula rows' residual variance.} A small $X$ means the common factors
        have already taken the systematic copula variation and the extra block
        will not buy much.
\end{itemize}

\paragraph{Acceptance.} The target remains the nine \code{wealth\_by\_income}
cells: $R^2 \approx 0.75$ on \code{mid40} and $\approx 0.8$ on \code{top10} is
what the truth's own predictability implies is reachable. The guardrail is that
the marginal and \code{income\_by\_wealth} cells, currently 92--99\%, must not
deteriorate --- they carry the paper's other results.

\end{document}
